\documentclass[paper=a4paper, fontsize=10pt, head_space=10mm, foot_space=17mm, gutter=17mm, line_length=185mm, twoside]{jlreq}
\usepackage{minted}
\usepackage{hyperref}
\pagestyle{plain}

\title{sshmount ソースリスト}
\author{美都}
\date{\today}

\begin{document}
\maketitle
\tableofcontents
\clearpage

\section{メインモジュール　main.rs}
\inputminted[linenos, breaklines]{rust}{src/main.rs}
\clearpage

\section{コマンドラインオプションの解析関係}
\subsection{コマンドラインオプションの定義 cmdline\_opt.rs}
\inputminted[linenos, breaklines]{rust}{src/cmdline_opt.rs}
\clearpage

\subsection{コマンドライン　remotehostの解析部 cmdline\_opt/remote\_host.rs}
\inputminted[linenos, breaklines]{rust}{src/cmdline_opt/remote_host.rs}
\clearpage

\section{ssh2 ログイン処理モジュール ssh\_connect.rs}
\inputminted[linenos, breaklines]{rust}{src/ssh_connect.rs}
\clearpage

\section{FUSE接続オプション生成モジュール fuse\_util.rs}
\inputminted[linenos, breaklines]{rust}{src/fuse_util.rs}
\clearpage

\section{FUSE　ファイルシステム本体と関連モジュール}
\subsection{ファイルシステムモジュール ssh\_filesystem.rs}
\inputminted[linenos, breaklines]{rust}{src/ssh_filesystem.rs}
\clearpage

\subsection{ファイルハンドル管理モジュール ssh\_filesystem/file\_handle.rs}
\inputminted[linenos, breaklines]{rust}{src/ssh_filesystem/file_handle.rs}
\clearpage

\subsection{Inode管理モジュール ssh\_filesystem/inode.rs}
\inputminted[linenos, breaklines]{rust}{src/ssh_filesystem/inode.rs}
\clearpage

\subsection{双方向ハッシュマップモジュール ssh\_filesystem/bi\_hash\_map.rs}
\inputminted[linenos, breaklines]{rust}{src/ssh_filesystem/bi_hash_map.rs}
\clearpage

\subsection{Mutex用拡張モジュール ssh\_filesystem/mutex\_lock\_ext.rs}
\inputminted[linenos, breaklines]{rust}{src/ssh_filesystem/mutex_lock_ext.rs}
\clearpage

\section{統合テスト用モジュール関係}
\subsection{単純なファイルアクセスのテストモジュール tests/simple\_file\_access.rs}
\inputminted[linenos, breaklines]{rust}{tests/simple_file_access.rs}
\clearpage

\subsection{共通ユーティリティー tests/common/mod.rs}
\inputminted[linenos, breaklines]{rust}{tests/common/mod.rs}
\clearpage

\subsection{テスト用サーバー設定のコンフィグファイルサンプル tests/test.env.sample}
 これは、cargo make testで使用する環境変数です。sshmountの接続設定を記述します。
 実際に使用するときは、リモートパスとマウントポイントを記述し、'test.env'にリネームします。
\inputminted[linenos, breaklines]{text}{tests/test.env.sample}

\subsection{主にテスト用に設定したcargo-makeのmakeファイル /Makefile.toml}
\inputminted[linenos, breaklines]{toml}{root/Makefile.toml}

\end{document}
